% *********************************************************************
% © 2016–2025 Jeremy Sylvestre
%
% Permission is granted to copy, distribute and/or modify this document
% under the terms of the GNU Free Documentation License, Version 1.3 or
% any later version published by the Free Software Foundation; with no
% Invariant Sections, no Front-Cover Texts, and no Back-Cover Texts. A
% copy of the license is included in the appendix entitled “GNU Free
% Documentation License” that appears in the output document of this
% PreTeXt source code. All trademarks™ are the registered® marks of
% their respective owners.
%
% *********************************************************************
\tdplotsetmaincoords{110}{30}
\begin{tikzpicture}[
	point/.style={circle,draw,very thin,fill,inner sep=0pt,minimum size=4pt},
	vector/.style={-latex},
	tdplot_main_coords
]

	% plane through the origin with normal vector (1,-1,2)
	% orthog basis used to draw sides: (1,1,0) and (1,-1,-1)

	% origin
	\node[point] at (0,0,0) (o) [label={[yshift=0.1cm]left:{$\zerovec$}}] {};
	% basis vectors for plane
	\coordinate (e1) at (1,1,0);
	\coordinate (e2) at (1,-1,-1);
	% but let's make the visualized e_2 a little longer
	\coordinate (e2_vis) at (1.5,-1.5,-1.5);
	% plane corners
	% sides of plane are vectors (4,4,0) and (3,-3,-3)
	\coordinate (P1) at (-2,0,1);
	\coordinate (P2) at (2,4,1);
	\coordinate (P3) at (5,1,-2);
	\coordinate (P4) at (1,-3,-2);
	% normal vector for plane
	\coordinate (n) at (1,-1,2);
	% perp space line endpoints
	%   (-3/4) n
	\coordinate (neg_n) at (-0.75,0.75,-1.5);
	%   (-1/2) n
	\coordinate (neg_n_plane) at (-0.375,0.375,-0.75);
	% external vector
	% v = (3/2) e_1 + (3/2) e_2 + (4/5) n
	% --> parallel to (38, -8, 1)
	\coordinate (v) at (3.8,-0.8,0.1);
	% proj v onto plane, equal to (3/2) e_1 + (3/2) e_2
	\coordinate (v_plane) at (3,0,-1.5);
	% proj v onto normal vector, equal to (4/5) n
	\coordinate (v_perp) at (0.8,-0.8,1.6);

	% fill/border/label for plane
	\filldraw[opacity=0.25,gray] (P1) to (P2) to (P3) to (P4) to cycle;
	\draw[dotted]  (P1) to (P2) to (P3) to (P4) to cycle;
	\node[xshift=-0.5cm,yshift=0.25cm] at (P3) {$U_2$};
	% U_2 = \Span \{\uvec{e}_1, \uvec{e}_2\}

	% perp space
	\draw (o) to (n) node[right] {$\orthogcmp{U}_2$};
	\draw[dashed] (o) -- (neg_n_plane);
	\draw (neg_n) -- (neg_n_plane);

	% plane vectors
	\draw[vector,thick] (o) to node[right,at end] {$\uvec{e}_1$} (e1);
	\draw[vector,thick] (o) to node[left,at end] {$\uvec{e}_2$} (e2_vis);
	% perp symbol: 20% in each direction
	\draw[thin] (0.2,0.2,0) to (0.4,0,-0.2) to (0.2,-0.2,-0.2);
	\draw[vector,thick] (o) to node[left,near end] {$\uvec{e}_3$} (v_perp);
	% \uvec{e}_3 = \proj_{\orthogcmp{U}_2} \uvec{v}_3 ;

	% third vector in original basis
	\draw[vector,thick] (o) to (v) node[above] {$\uvec{v}_3$};
	\draw[dashed] (v_perp) to (v) to (v_plane);
	\draw[vector,thick] (0.36,0,-0.18) to (v_plane); % don't draw right from (o) because of e1/e2 perp symbol
	\node at (1,0,-2) (v3_label) {$\proj_{U_2} \uvec{v}_3$};
	\draw[densely dotted] (v3_label.north) to (1.5,0,-0.75);
	% perp symbol for projection of v onto plane:
	% (1/8) of v_perp, 10% of -(v_plane)
	\draw[thin] (3.1, -0.1, -1.3) to (2.8,-0.1,-1.15) to (2.7,0,-1.35);

\end{tikzpicture}
